\begin{abstract}
We disprove the scalar-curvature extension conjecture in its unrestricted
all-dimensions form. In sufficiently high dimension, we construct a Ricci
flow on a closed manifold whose scalar curvature remains uniformly bounded
while full curvature diverges at a finite maximal time. In one fixed
sufficiently high dimension, the examples have two-sided power-law curvature
blowup with arbitrarily large exponents.
\end{abstract}

\section{Introduction}

A Ricci flow on a smooth manifold is a family of Riemannian metrics
\(g(t)\) satisfying \(\partial_tg=-2\Ric_g\). Here a closed manifold
is compact and has no boundary, \(\Rm\) denotes its full curvature
tensor, and \(R=\operatorname{tr}_g\Ric\) denotes scalar curvature.
A smooth Ricci flow on a closed manifold can fail to extend at a finite time
only if its full curvature becomes unbounded~\cite{Hamilton1995}.
The scalar-curvature extension conjecture asks whether scalar curvature
alone must detect such a singularity; see~\cite[p.~756]{Bamler2018}.
A uniform bound for the Ricci tensor is sufficient for extension
by \v{S}e\v{s}um's Theorem~\cite{Sesum2005}, but the trace of that tensor
contains substantially less information. The question concerns smooth
extension on the original manifold; continuation through a singular space
is a different conclusion.

Bounded scalar curvature imposes a necessary restriction on the rate of
full-curvature blowup. Wang proved that, at a finite singular time \(T\),
the product of \(T-t\), the square root of the maximum full curvature,
and the square root of the maximum absolute scalar curvature has a
strictly positive upper limit~\cite[Theorem~3]{Wang2012}.
Consequently, a closed flow with bounded scalar curvature must satisfy
\[
\limsup_{t\uparrow T}(T-t)^2\max_M|\Rm|(\cdot,t)>0.
\]
This is a subsequential restriction; it does not assert a lower bound at
every sufficiently late time. It explains why examples with curvature
concentration faster than the parabolic rate are relevant to the
extension question.

Symmetry has made it possible to construct and study singularities with
several distinct geometric scales. Angenent and Knopf constructed
rotationally symmetric neckpinches on spheres~\cite{AK2004}.
A finite-time singularity is of
\emph{Type II} when \((T-t)\max_M|\Rm|\) is unbounded.
Gu and Zhu constructed rotationally symmetric Type-II singularities on
spheres~\cite{GuZhu2008}. Angenent, Isenberg, and Knopf later constructed
degenerate neckpinches with prescribed rates, a Bryant-soliton tip model,
and a shrinking-cylinder parabolic model~\cite{AIK2015}.
The Type-II constructions establish that full curvature can concentrate faster
than the parabolic scale, but do not supply the scalar bound in
Theorem~\ref{thm:main}.

Stolarski constructed closed, doubly warped Ricci flows that develop
conical singularities with arbitrarily fast curvature blow-up and approach
a Ricci-flat cone at the parabolic scale~\cite{Stolarski2019}.
Stolarski's discussion identifies bounded scalar curvature as a possible
feature of these examples. The construction does not supply the
scalar-curvature bound proved below.
We use its quantitative profile estimates and establish that bound for a
choice of sufficiently large dimension and mode index.
The geometric cap models considered in his formal discussion belong to
the cohomogeneity-one Ricci-flat setting developed by
B\"ohm~\cite{Bohm1999,Stolarski2019}. Our cap estimate uses only
the quantitative construction inputs stated in
Proposition~\ref{prop:input}; it does not require convergence to a
prescribed complete cap metric.

\begin{theorem}\label{thm:main}
There are an integer \(q\ge10\), a time \(T>0\), and a smooth Ricci flow
\(g(t)\), \(0\le t<T\), on the closed connected manifold
\(M=\Sp^2\times\Sp^{q+1}\), with smooth initial metric, such that
\[
\sup_{M\times[0,T)}|R_{g(t)}|<\infty,
\qquad
\lim_{t\uparrow T}\max_M|\Rm_{g(t)}|=\infty.
\]
In particular, \(T\) is the finite maximal existence time of this flow.
\end{theorem}

Theorem~\ref{thm:main} refutes this conjecture in its unrestricted
all-dimensions formulation: the assertion that every finite-time
singularity of a closed Ricci flow in dimension at least four has unbounded
scalar curvature. The example is in sufficiently high dimension; no
dimension-four conclusion is asserted.

The estimates also determine the full-curvature rate of the examples.
Corollary~\ref{cor:rates} gives two-sided power-law bounds with an unbounded
discrete set of exponents, while the dimension remains fixed. Thus the
scalar bound is compatible with arbitrarily fast power-law Type-II
singularities, and the upper rate is controlled as well as the lower rate.

\subsection*{The two scales and the proof}

The examples have two sphere factors whose radii vary along an interval.
Off the two pole orbits the metric is
\[
ds^2+\phi(s,t)^2g_{\Sp^2}+r(s,t)^2g_{\Sp^q},
\]
where \(s\) is arclength at the indicated time and the sphere metrics
have sectional curvature one. At each endpoint the \(\Sp^q\) factor
collapses smoothly, leaving a positive-radius \(\Sp^2\) orbit.
The singularity develops when that remaining orbit also shrinks.

Write \(\delta=T-t\). The parabolic length scale is \(\sqrt\delta\),
where the rescaled metric closely approximates a Ricci-flat cone.
The smaller cap scale is \(\theta=\delta^\sigma\), with a fixed
exponent \(\sigma>1/2\) supplied by the construction. These two
scales have different roles: the cone approximation gives small Ricci
curvature on fixed parabolic annuli, whereas smoothness at the pole must
be controlled on the cap scale. A small Ricci tensor at the larger scale
does not by itself control the smaller cap.

The additional estimates have two features. First, Stolarski's strict
one-sided comparison with the cone and monotonicity of a logarithmic
slope, together with the ordinary scalar lower bound, control the size of
the cap. We adapt the scalar-sign obstruction in his
Proposition~5.3~\cite{Stolarski2019} to obtain this quantitative estimate.
Curvature
point-picking then gives a full-curvature bound without assuming convergence
to a cap model. Second, on radial cylinders moving with the dominant drift,
one-dimensional parabolic estimates give a Ricci reaction bound with leading
coefficient \(2q\). The corresponding diffusion coefficient leaves a positive
margin in large dimension. This yields
\( |\Ric|\le Cr^{-e}\) for an exponent \(0<e<1\), where \(r\) is
the radius of the \(\Sp^q\) factor in the doubly warped metric.
The square of this bound admits a bounded scalar-curvature supersolution.
These are the new estimates beyond the imported singular-flow construction.

Section~\ref{sec:construction} states the precise construction input and
translates its scales. Section~\ref{sec:geometry} obtains the warping
identities, the strict cone gap, and the parabolic-annulus estimates.
Section~\ref{sec:cap} proves the cap and radius curvature bounds by two
curvature-record arguments and Shi's derivative estimates.
Section~\ref{sec:reaction} proves the dimension-explicit reaction bound;
the analytic estimates take place in one space dimension, which is what
keeps their constants independent of \(q\).
Section~\ref{sec:barriers} fixes the parameters in order and completes
the Ricci and scalar comparisons, then records the resulting curvature rates.

\paragraph{Conventions.}
All tensor norms and geometric differential operators are taken with respect
to \(g(t)\), and \(\dt g=-2\Ric\).
The unit sphere has sectional curvature one. A constant \(C\) may change
between occurrences. Unless explicitly stated otherwise, it may depend on
all fixed parameters of the chosen flow, but never on \(t\uparrow T\).
Constants used to choose the dimension will be identified as independent
of both the dimension parameter \(q\) and the mode index \(k\).
